\documentclass[11pt,a4paper]{article}

\usepackage[utf8]{inputenc}
\usepackage[T1]{fontenc}
\usepackage{amsmath,amssymb}
\usepackage{graphicx}
\usepackage{hyperref}
\usepackage{xcolor}
\usepackage[margin=1in]{geometry}

\renewcommand{\baselinestretch}{2.0}

\hypersetup{
    colorlinks=true,
    linkcolor=blue,
    citecolor=blue,
    urlcolor=blue
}

\title{Implicit Geographic Inference in LLM Medical Triage:\\
Language-Driven Disparities in Emergency Recommendations}

\author{
    Qi Han Wong, B.A.\\
    Independent Researcher\\
    \texttt{wongqihan@gmail.com}
}

\date{}

\begin{document}

\maketitle

% ============================================================================
\begin{abstract}

\noindent\textbf{Background:}
Large language models are increasingly used in healthcare-adjacent applications, including symptom checkers and triage assistants. Whether these systems produce equitable recommendations across languages is not well understood.

\noindent\textbf{Methods:}
We presented identical neurological symptoms (persistent headache, blurred vision, nausea in a 38-year-old patient) to Gemini 3.5 Flash in eight languages (English, Spanish, Chinese, Hindi, Japanese, Arabic, French, Russian) with 30 runs per condition ($n = 570$ total API calls). The primary outcome was the rate of emergency room (ER) recommendations. We conducted location-anchoring experiments and a back-translation control to isolate the mechanism. We report 95\% Wilson score confidence intervals for all proportions.

\noindent\textbf{Results:}
ER recommendation rates ranged from 0\% (Japanese, Hindi) to 33.3\% (French) despite nearly identical severity scores (7.7--8.0 out of 10) across all languages. Adding a single sentence specifying a US location increased ER recommendations by up to 76.7 percentage points for non-English prompts. The reverse anchor (English prompt with Tokyo location) reduced the ER rate from 30.0\% to 6.7\%. A back-translation control (Japanese to English) produced ER rates comparable to the English baseline (36.7\% vs 30.0\%), confirming that the disparity reflects implicit geographic inference rather than translation quality differences.

\noindent\textbf{Conclusions:}
The model uses input language as a proxy for the patient's geographic location and applies region-specific healthcare norms to its triage recommendations. This behavior produces appropriate recommendations for monolingual users in their home country but may disadvantage multilingual users, immigrants, and expatriates whose language does not match their geographic location. Explicitly specifying the patient's location in the system prompt eliminates the disparity.

\end{abstract}

% ============================================================================
\section{Introduction}

Large language models are increasingly deployed in healthcare-adjacent applications, including symptom checkers, triage assistants, and patient-facing chatbots.$^{1,2}$ A critical assumption underlying these deployments is that clinical recommendations are driven by the medical content of the input rather than superficial features such as the language in which the query is written.

Prior work has documented racial and gender biases in LLM-generated medical content. Zack et al.$^3$ found that GPT-4 perpetuates race-based clinical decision-making when racial identifiers are included in prompts. Omiye et al.$^4$ showed that large language models propagate race-based medicine by reproducing debunked claims about biological differences between racial groups. Cross-lingual performance gaps are also well-documented: models consistently perform worse on non-English benchmarks,$^{5,6}$ and safety alignment degrades in low-resource languages.$^{7,8}$ However, most prior work focuses on \textit{capability} gaps (the model understands less) rather than \textit{normative} shifts (the model applies different standards). LLMs also encode geographic and cultural knowledge from their training corpora,$^{9,10}$ but whether this internalized knowledge actively modifies clinical recommendations has not been tested.

We test whether a production LLM produces different medical triage recommendations for identical symptoms based solely on the language of the patient prompt. We present a single neurological symptom profile to Gemini 3.5 Flash in eight languages and measure the recommended clinical action. Our key finding is that the model uses language as a proxy for geographic location, then applies the healthcare norms of the inferred country. The model assigns nearly identical severity scores across all languages, confirming equivalent clinical comprehension, but recommends emergency room visits at rates ranging from 0\% (Japanese, Hindi) to 33.3\% (French).

We validate this mechanism through three controls: (1) adding a US location anchor to non-English prompts increases ER rates by up to 76.7 percentage points; (2) adding a Tokyo location anchor to the English prompt reduces the ER rate from 30.0\% to 6.7\%; and (3) back-translating the Japanese prompt to English produces ER rates comparable to the English baseline, ruling out translation quality as a confound. The complete experiment code, prompts, and raw results are publicly available.\footnote{\url{https://github.com/wongqihan/ai-behavioral-experiments}}

% ============================================================================
\section{Methods}

\subsection{Model and Configuration}

We used Gemini 3.5 Flash (Google, May 2026) accessed via the public API with temperature 0.3 and 30 runs per condition. We chose a low temperature to maximize response consistency while preserving stochastic variation for statistical analysis.

\subsection{Symptom Profile}

We used a single neurological symptom profile describing a 38-year-old patient with persistent headache for two weeks unresponsive to painkillers, progressive blurred vision (especially when reading or using screens), morning nausea, occasional visual disturbances (spots), and no prior history of similar symptoms. This profile was selected because it is clinically ambiguous---consistent with conditions ranging from tension headache (low urgency) to increased intracranial pressure (high urgency)---creating a decision boundary where contextual priors can influence the recommended action.

\subsection{Prompt Design}

All symptom prompts were manually authored (not machine-translated) to match native phrasing conventions in each language: English, Spanish, Chinese (Simplified), Hindi, Japanese, Arabic, French, and Russian. These eight languages were selected to represent diverse writing systems, geographic regions, and healthcare infrastructure norms.

The system prompt enforced structured JSON output with fields for diagnosis (English), severity (1--10), urgency (Emergency/Urgent/Routine), action (ER/Doctor appointment/Self-care), recommended tests (list), and reasoning (English). The system prompt was always in English regardless of user prompt language.

\subsection{Test Conditions}

We conducted four test conditions: (1) \textbf{Baseline:} symptom prompts in 8 languages with no location context ($8 \times 30 = 240$ calls); (2) \textbf{US Location Anchor:} same prompts with an appended sentence in each language specifying a US location ($8 \times 30 = 240$ calls); (3) \textbf{Reverse Anchor:} English prompt with ``Assume the patient is located in Tokyo, Japan'' or ``Assume the patient is located in Mumbai, India'' ($2 \times 30 = 60$ calls); and (4) \textbf{Back-translation Control:} the Japanese prompt translated to English by the same model, then triaged ($1 \times 30 = 30$ calls). Total: 570 API calls.

\subsection{Statistical Analysis}

The primary outcome was the ER recommendation rate (proportion of runs where action = ``ER''). Secondary outcomes included average severity score and urgency classification. All API calls returned valid, parseable JSON responses (0 parse failures). We report 95\% Wilson score confidence intervals for all proportions.

\subsection{AI Disclosure}

AI tools (Anthropic Claude and Google Gemini) were used to assist with manuscript preparation, including drafting and editing text. All content was reviewed, verified, and approved by the author. The experiment code was written by the author with AI coding assistance.

% ============================================================================
\section{Results}

\subsection{Baseline: Language Shifts ER Rates}

Table~\ref{tab:baseline} shows ER recommendation rates across eight languages with no location context. Severity scores were nearly identical (7.7--8.0), confirming equivalent clinical assessment regardless of language. The divergence was entirely in the recommended action. French (33.3\%) and English/Arabic (30.0\%) had the highest ER rates, while Hindi and Japanese had 0\% ER rates. The confidence intervals for the high-ER group and the zero-ER group did not overlap.

\begin{table}[h]
\centering
\caption{Baseline ER recommendation rates by language (no location context, $n=30$ per language).}
\label{tab:baseline}
\begin{tabular}{lcccc}
\hline
Language & ER Count & ER \% & 95\% CI & Avg Severity \\
\hline
French   & 10/30 & 33.3\% & [19.2\%, 51.2\%] & 7.7 \\
English  & 9/30  & 30.0\% & [16.7\%, 47.9\%] & 7.7 \\
Arabic   & 9/30  & 30.0\% & [16.7\%, 47.9\%] & 8.0 \\
Russian  & 7/30  & 23.3\% & [11.8\%, 40.9\%] & 8.0 \\
Chinese  & 6/30  & 20.0\% & [9.5\%, 37.3\%]  & 8.0 \\
Spanish  & 4/30  & 13.3\% & [5.3\%, 29.7\%]  & 7.8 \\
Hindi    & 0/30  & 0.0\%  & [0.0\%, 11.4\%]  & 8.0 \\
Japanese & 0/30  & 0.0\%  & [0.0\%, 11.4\%]  & 8.0 \\
\hline
\end{tabular}
\end{table}

\subsection{US Location Anchor}

Adding a single sentence specifying a US location caused ER rates to surge across all non-English languages (Table~\ref{tab:anchor}). English showed minimal change (+10.0 pp), consistent with the model already assuming a US location for English queries. The largest shifts occurred for Chinese (+76.7 pp) and Hindi (+73.3 pp). The English shift has heavily overlapping confidence intervals with the English baseline and is not statistically significant, while all non-English shifts produce non-overlapping intervals with their respective baselines.

\begin{table}[h]
\centering
\caption{Effect of US location anchor on ER recommendation rates ($n=30$ per condition).}
\label{tab:anchor}
\begin{tabular}{lcccc}
\hline
Language & Default ER\% & + US Anchor ER\% & Shift & Anchor 95\% CI \\
\hline
Chinese  & 20.0\% & 96.7\% & +76.7 pp & [83.3\%, 99.4\%] \\
Hindi    & 0.0\%  & 73.3\% & +73.3 pp & [55.6\%, 85.8\%] \\
Arabic   & 30.0\% & 90.0\% & +60.0 pp & [74.4\%, 96.5\%] \\
Spanish  & 13.3\% & 70.0\% & +56.7 pp & [52.1\%, 83.3\%] \\
Russian  & 23.3\% & 76.7\% & +53.3 pp & [59.1\%, 88.2\%] \\
Japanese & 0.0\%  & 46.7\% & +46.7 pp & [30.2\%, 63.9\%] \\
French   & 33.3\% & 53.3\% & +20.0 pp & [36.1\%, 69.8\%] \\
English  & 30.0\% & 40.0\% & +10.0 pp & [24.6\%, 57.7\%] \\
\hline
\end{tabular}
\end{table}

\subsection{Reverse Anchor and Back-Translation}

Applying a foreign location anchor to an English prompt reversed the effect (Table~\ref{tab:controls}). English with a Tokyo anchor produced an ER rate of 6.7\%, and English with a Mumbai anchor produced 0.0\%, both substantially below the English baseline of 30.0\%. This confirms that language alone does not drive the disparity; explicitly overriding the inferred location changes the recommended action.

The back-translation control (Japanese prompt translated to English by the same model) produced an ER rate of 36.7\%, with a confidence interval [21.9\%, 54.5\%] overlapping almost entirely with the English baseline [16.7\%, 47.9\%]. The 0\% ER rate for the Japanese prompt is not caused by the model failing to understand the symptoms. It is caused by the model inferring a Japanese location and applying Japanese healthcare routing norms.

\begin{table}[h]
\centering
\caption{Reverse anchor and back-translation controls ($n=30$ per condition).}
\label{tab:controls}
\begin{tabular}{lccc}
\hline
Condition & ER Count & ER \% & 95\% CI \\
\hline
English (baseline)            & 9/30 & 30.0\% & [16.7\%, 47.9\%] \\
English + Tokyo anchor        & 2/30 & 6.7\%  & [1.8\%, 21.3\%]  \\
English + Mumbai anchor       & 0/30 & 0.0\%  & [0.0\%, 11.4\%]  \\
Japanese (original)           & 0/30 & 0.0\%  & [0.0\%, 11.4\%]  \\
Japanese $\rightarrow$ English & 11/30 & 36.7\% & [21.9\%, 54.5\%] \\
\hline
\end{tabular}
\end{table}

% ============================================================================
\section{Discussion}

Our results demonstrate that a production large language model uses input language as a proxy for the patient's geographic location and applies region-specific healthcare norms to its triage recommendations. The inference chain is:

\begin{center}
\textit{Input language} $\rightarrow$ \textit{Inferred country} $\rightarrow$ \textit{Regional healthcare norms} $\rightarrow$ \textit{Recommended action}
\end{center}

Japanese prompts are routed through Japan's clinic-first pathway; English and French prompts are routed through US and French emergency medicine norms; Hindi prompts are routed through India's conservative triage approach. Notably, the model does not mention geography in its reasoning field. The reasoning text across all languages describes the same clinical picture (intracranial pressure risk, papilledema, need for MRI). The geographic inference is implicit---it shapes the recommended action without surfacing in the model's stated rationale.

This behavior is appropriate for a monolingual user located in the country associated with their language. A Japanese speaker in Tokyo \textit{should} receive recommendations aligned with Japanese healthcare infrastructure. The failure mode arises for users whose language does not match their location: a Hindi-speaking immigrant in San Francisco receives 0\% ER recommendations for symptoms that generate 30\% ER recommendations in English; multilingual users receive different recommendations depending on which language they choose to describe their symptoms.

The most direct mitigation is to explicitly anchor the patient's geographic location in the system prompt. Our US anchor results demonstrate that the model already knows how to apply location-appropriate norms when given explicit context. The fix is not to remove geographic awareness but to decouple it from language. When the patient's location is unknown, the system should ask rather than infer.

\subsection{Limitations}

This study has several limitations. We tested one model (Gemini 3.5 Flash). Preliminary testing of Claude (Anthropic) and GPT (OpenAI) models did not replicate this pattern, with both defaulting to emergency recommendations regardless of language, suggesting the geographic inference behavior may reflect model-specific localization training rather than a universal LLM property. We tested one symptom profile; different clinical presentations may show different degrees of geographic sensitivity. Prompts were manually authored but not independently verified by native-speaking medical professionals. We did not compare to human clinician recommendations from each country. At $n=30$ per condition, some pairwise comparisons have overlapping confidence intervals. The system prompt was always in English, which may introduce asymmetry for non-English user prompts.

% ============================================================================
\section{Conclusion}

A production large language model produces systematically different medical triage recommendations based on prompt language, even when clinical content is identical. The mechanism is implicit geographic inference: the model uses language as a proxy for the patient's location and applies region-specific healthcare norms. For developers deploying LLMs in medical triage, our findings suggest a concrete mitigation: decouple geographic context from language by explicitly specifying the patient's location. Language should determine \textit{how} the model communicates, not \textit{what} it recommends.

Code, data, and full results are available at:\\
\url{https://github.com/wongqihan/ai-behavioral-experiments}

% ============================================================================
\section*{References}

\begin{enumerate}
\item Singhal K, Azizi S, Tu T, et al. Large language models encode clinical knowledge. Nature 2023;620(7972):172-180. DOI: 10.1038/s41586-023-06291-2.
\item Nori H, King N, McKinney SM, Carignan D, Horvitz E. Capabilities of GPT-4 on medical challenge problems. arXiv preprint arXiv:2303.13375. 2023.
\item Zack T, Lehman E, Suzgun M, et al. Assessing the potential of GPT-4 to perpetuate racial and gender biases in health care: a model evaluation study. Lancet Digit Health 2024;6(1):e12-e22. DOI: 10.1016/S2589-7500(23)00225-X.
\item Omiye JA, Lester JC, Spichak S, Rotemberg V, Daneshjou R. Large language models propagate race-based medicine. NPJ Digit Med 2023;6(1):195. DOI: 10.1038/s41746-023-00939-z.
\item Ahuja K, Diddee H, Hada R, et al. MEGA: Multilingual evaluation of generative AI. In: Proceedings of the 2023 Conference on Empirical Methods in Natural Language Processing. 2023.
\item Lai VD, Ngo NT, Veyseh APB, et al. ChatGPT beyond English: towards a comprehensive evaluation of large language models in multilingual learning. In: Findings of EMNLP. 2023.
\item Deng Y, Zhang W, Pan SJ, Bing L. Multilingual jailbreak challenges in large language models. In: Proceedings of ICLR. 2024.
\item Yong ZX, Menghini C, Bach SH. Low-resource languages jailbreak GPT-4. arXiv preprint arXiv:2310.02446. 2024.
\item Manvi R, Khanna S, Burke M, Lobell D, Ermon S. Large language models are geographically biased. In: Proceedings of ICML. 2024.
\item Li H, Jiang L, Hwang JD, et al. Culture-Gen: revealing global cultural perception in language models through natural language prompting. In: Proceedings of COLM. 2024.
\item Chen IY, Pierson E, Rose S, et al. Ethical machine learning in healthcare. Annu Rev Biomed Data Sci 2021;4:123-144. DOI: 10.1146/annurev-biodatasci-092820-114757.
\item Obermeyer Z, Powers B, Vogeli C, Mullainathan S. Dissecting racial bias in an algorithm used to manage the health of populations. Science 2019;366(6464):447-453. DOI: 10.1126/science.aax2342.
\end{enumerate}

\end{document}
